\documentclass[conference]{IEEEtran}

\usepackage{amsmath,amssymb,amsthm}
\usepackage{booktabs}
\usepackage{array}
\usepackage{cite}
\usepackage{hyperref}
\usepackage{graphicx}
\usepackage{xcolor}

\theoremstyle{definition}
\newtheorem{definition}{Definition}[section]
\newtheorem{theorem}{Theorem}[section]
\newtheorem{proposition}{Proposition}[section]

\begin{document}

\title{Classical Time-Series Models: A Comprehensive Survey\\
{\large Data Science --- Advanced Predictive Models}}

\author{\IEEEauthorblockN{Surya Rao Rayarao}
\IEEEauthorblockA{\textit{Department of Statistics and Data Sciences} \\
\textit{Department of Computer Science} \\
\textit{The University of Texas at Austin}\\
Austin, TX, USA \\
suryarao.r@utexas.edu}
\and
\IEEEauthorblockN{Naga Donikena}
\IEEEauthorblockA{\textit{Department of Statistics and Data Sciences} \\
\textit{Department of Computer Science} \\
\textit{The University of Texas at Austin}\\
Austin, TX, USA \\
naga.donikena@utexas.edu}
}

\maketitle

%=============================================================================
\begin{abstract}
Classical time-series models form the statistical backbone of temporal data
analysis, underpinning applications from economic forecasting to sensor
monitoring. This survey provides a rigorous treatment of the foundational
model families: autoregressive (AR), moving-average (MA), their integrated
(ARIMA) and seasonal (SARIMA) extensions, exponential smoothing state-space
models, and error-correction models for co-integrated series. We develop the
mathematical formalism of stationarity, autocorrelation structure, and
information criteria–driven model selection, then describe estimation
algorithms including the Yule--Walker equations, conditional and exact maximum
likelihood, and the Kalman filter recursion. Worked numerical examples
illustrate identification, estimation, diagnostic checking, and multi-step
forecasting for canonical datasets. We survey the deployment of these models
in demand forecasting, financial risk modeling, epidemiological surveillance,
and industrial process control. We also analyze the assumptions and
limitations---linearity, Gaussianity, fixed structure---that motivate
connections to ARCH/GARCH models, state-space and structural time-series
models, and modern machine-learning approaches. The survey is intended for
graduate students in data science seeking a unified, mathematically grounded
reference for dependent-data modeling and forecasting.
\end{abstract}

%=============================================================================
\section{Introduction}
\label{sec:intro}

Time-ordered observations arise in virtually every quantitative discipline.
Unlike independent samples, successive measurements of a stochastic process
typically exhibit \emph{serial dependence}: knowing yesterday's stock price,
temperature, or patient heart rate provides information about today's value.
Ignoring this dependence leads to invalid confidence intervals, sub-optimal
forecasts, and missed signals about latent dynamics. Classical time-series
analysis provides a principled mathematical framework for modeling, estimating,
diagnosing, and forecasting serially dependent data.

The intellectual lineage is long. \emph{Periodogram} analysis traces to
Schuster~\cite{schuster1898investigation} and Yule~\cite{yule1927periodicity},
who introduced autoregressive (AR) models to explain quasi-periodic behavior
in sunspot data. Slutzky~\cite{slutzky1927summation} showed that summation of
random shocks produces moving-average (MA) processes with periodic-looking
realizations. Wold's decomposition theorem~\cite{wold1938study} unified these
ideas, proving that every covariance-stationary process admits an infinite MA
representation. Box and Jenkins~\cite{box2015time} operationalized these
theoretical results into the iterative \emph{identification--estimation--%
diagnostic} cycle that remains the standard applied workflow. Holt--Winters
exponential smoothing~\cite{holt2004forecasting,winters1960forecasting}
provided a computationally simpler alternative that proved especially popular
in operations research and inventory management.

This survey synthesizes these classical contributions from a unified perspective
aligned with four core themes of advanced predictive modeling:
\begin{enumerate}
  \item \textbf{Identifying dependencies in structured data} --- autocorrelation
        and partial autocorrelation functions, unit-root and co-integration
        tests.
  \item \textbf{Applying appropriate models for dependent data} --- AR, MA,
        ARMA, ARIMA, SARIMA, exponential smoothing, and error-correction models.
  \item \textbf{Making forecasts accounting for dependence and structure} ---
        minimum mean-squared-error (MMSE) predictors, prediction intervals, and
        multi-step horizon strategies.
  \item \textbf{Identifying latent structure in complex data} --- trend,
        seasonality, and irregular components; state-space decompositions.
\end{enumerate}

\subsection*{Paper Organization}
Section~\ref{sec:prelim} establishes mathematical notation and foundational
concepts. Section~\ref{sec:theory} develops the core theory of AR, MA, ARMA,
integrated, and seasonal models. Section~\ref{sec:algorithms} describes
estimation and model-selection algorithms. Section~\ref{sec:examples} works
through numerical examples. Section~\ref{sec:applications} covers data-science
applications. Section~\ref{sec:limits} discusses assumptions and limitations.
Section~\ref{sec:related} connects classical models to related methods.
Section~\ref{sec:conclusion} concludes.

%=============================================================================
\section{Mathematical Preliminaries}
\label{sec:prelim}

\subsection{Notation and Conventions}

Let $\{Y_t\}_{t \in \mathbb{Z}}$ denote a real-valued discrete-time stochastic
process indexed by integers. We observe the finite realization
$\mathbf{y} = (y_1, y_2, \ldots, y_T)^{\top}$.

The \emph{backshift operator} $B$ is defined by $BY_t = Y_{t-1}$ and
$B^k Y_t = Y_{t-k}$ for any integer $k$. The \emph{difference operator} is
$\nabla = 1 - B$, so $\nabla Y_t = Y_t - Y_{t-1}$.

\begin{definition}[Strict Stationarity]
A process $\{Y_t\}$ is \emph{strictly stationary} if for every finite set of
indices $\{t_1,\ldots,t_k\}$ and every integer $h$,
\[
  (Y_{t_1},\ldots,Y_{t_k}) \stackrel{d}{=} (Y_{t_1+h},\ldots,Y_{t_k+h}).
\]
\end{definition}

\begin{definition}[Covariance Stationarity]
$\{Y_t\}$ is \emph{covariance-stationary} (weakly stationary) if
\begin{align}
  \mathbb{E}[Y_t] &= \mu \quad \forall\, t, \label{eq:mean_stat} \\
  \mathrm{Cov}(Y_t, Y_{t+h}) &= \gamma(h) \quad \forall\, t, h, \label{eq:cov_stat}
\end{align}
where both $\mu$ and $\gamma(h)$ are finite.
\end{definition}

\subsection{Autocovariance and Autocorrelation}

The \emph{autocovariance function} (ACVF) at lag $h$ is
$\gamma(h) = \mathrm{Cov}(Y_t, Y_{t+h})$. The \emph{autocorrelation function}
(ACF) is
\begin{equation}
  \rho(h) = \frac{\gamma(h)}{\gamma(0)}, \quad |\rho(h)| \leq 1.
  \label{eq:acf}
\end{equation}

The \emph{sample ACF} at lag $h$ is
\begin{equation}
  \hat{\rho}(h) = \frac{\sum_{t=1}^{T-h}(y_t - \bar{y})(y_{t+h} - \bar{y})}
                        {\sum_{t=1}^{T}(y_t - \bar{y})^2},
  \label{eq:sacf}
\end{equation}
where $\bar{y} = T^{-1}\sum_{t=1}^T y_t$.

The \emph{partial autocorrelation function} (PACF) $\phi_{hh}$ at lag $h$ is
the last coefficient in the best linear prediction of $Y_t$ by
$Y_{t-1},\ldots,Y_{t-h}$. Formally, if
$\hat{Y}_t^{(h)} = \phi_{h1}Y_{t-1} + \cdots + \phi_{hh}Y_{t-h}$
is the projection, then $\phi_{hh}$ solves the Yule--Walker system of order
$h$.

\subsection{Wold Decomposition}

\begin{theorem}[Wold, 1938~\cite{wold1938study}]
Every zero-mean covariance-stationary process $\{Y_t\}$ that is not
deterministic can be written as
\begin{equation}
  Y_t = \sum_{j=0}^{\infty} \psi_j \varepsilon_{t-j} + \kappa_t,
  \label{eq:wold}
\end{equation}
where $\{\varepsilon_t\}$ is white noise with variance $\sigma^2$,
$\psi_0 = 1$, $\sum_{j=0}^{\infty}\psi_j^2 < \infty$, and $\kappa_t$ is
perfectly predictable from its own past.
\end{theorem}

The Wold decomposition justifies restricting attention to ARMA models, which
provide parsimonious rational approximations to the infinite MA polynomial.

\subsection{White Noise and Innovation Sequences}

\begin{definition}[White Noise]
$\{\varepsilon_t\}$ is \emph{white noise} (WN$(0,\sigma^2)$) if
$\mathbb{E}[\varepsilon_t]=0$, $\mathrm{Var}(\varepsilon_t)=\sigma^2$, and
$\mathrm{Cov}(\varepsilon_t, \varepsilon_s)=0$ for $t \neq s$.
\end{definition}

When $\varepsilon_t \stackrel{\text{iid}}{\sim} \mathcal{N}(0,\sigma^2)$ we
speak of \emph{Gaussian white noise}, enabling exact likelihood computation.

%=============================================================================
\section{Core Theory}
\label{sec:theory}

\subsection{Autoregressive Models}

\begin{definition}[AR($p$) Process]
$\{Y_t\}$ follows an \emph{autoregressive process of order $p$} if
\begin{equation}
  Y_t = c + \phi_1 Y_{t-1} + \phi_2 Y_{t-2} + \cdots + \phi_p Y_{t-p}
        + \varepsilon_t,
  \label{eq:ar_p}
\end{equation}
where $c$ is an intercept, $\phi_1,\ldots,\phi_p$ are autoregressive
coefficients, and $\varepsilon_t \sim \mathrm{WN}(0,\sigma^2)$.
\end{definition}

Using the backshift operator, \eqref{eq:ar_p} is written compactly as
$\Phi(B)Y_t = c + \varepsilon_t$, where
$\Phi(z) = 1 - \phi_1 z - \phi_2 z^2 - \cdots - \phi_p z^p$
is the \emph{autoregressive (AR) polynomial}.

\begin{proposition}[Stationarity Condition]
The AR($p$) process is covariance-stationary if and only if all roots of
$\Phi(z) = 0$ lie \emph{outside} the unit circle in the complex plane,
i.e., $|z_i| > 1$ for all roots $z_i$.
\end{proposition}

For an AR(1) model $Y_t = \phi Y_{t-1} + \varepsilon_t$, stationarity requires
$|\phi| < 1$. Under stationarity:
\begin{align}
  \mu &= \frac{c}{1 - \phi_1 - \cdots - \phi_p}, \label{eq:ar_mean}\\
  \gamma(0) &= \sigma^2 / (1 - \phi_1 \rho(1) - \cdots - \phi_p \rho(p)),
  \label{eq:ar_var}
\end{align}
and the ACVF satisfies the \emph{Yule--Walker equations}:
\begin{equation}
  \gamma(h) = \phi_1\gamma(h-1) + \cdots + \phi_p\gamma(h-p), \quad h \geq 1.
  \label{eq:yw}
\end{equation}

The PACF of an AR($p$) process has a \emph{sharp cutoff} after lag $p$:
$\phi_{hh} = 0$ for $h > p$.

\subsection{Moving-Average Models}

\begin{definition}[MA($q$) Process]
$\{Y_t\}$ follows a \emph{moving-average process of order $q$} if
\begin{equation}
  Y_t = \mu + \varepsilon_t + \theta_1\varepsilon_{t-1} + \cdots
        + \theta_q\varepsilon_{t-q},
  \label{eq:ma_q}
\end{equation}
i.e., $Y_t = \mu + \Theta(B)\varepsilon_t$ with
$\Theta(z) = 1 + \theta_1 z + \cdots + \theta_q z^q$.
\end{definition}

Every finite-order MA process is covariance-stationary. The ACF satisfies:
\begin{equation}
  \rho(h) = \frac{\sum_{j=0}^{q-h}\theta_j\theta_{j+h}}
                 {\sum_{j=0}^{q}\theta_j^2}, \quad 1 \leq h \leq q,
  \label{eq:ma_acf}
\end{equation}
and $\rho(h) = 0$ for $|h| > q$ --- a \emph{sharp cutoff} at lag $q$.

\begin{definition}[Invertibility]
An MA($q$) process is \emph{invertible} if all roots of $\Theta(z) = 0$ lie
outside the unit circle. Invertibility ensures the unique infinite AR
representation $\varepsilon_t = \Theta^{-1}(B)(Y_t - \mu)$.
\end{definition}

\subsection{ARMA Models}

\begin{definition}[ARMA($p,q$) Process]
$\{Y_t\}$ follows an ARMA($p,q$) process if
\begin{equation}
  \Phi(B)(Y_t - \mu) = \Theta(B)\varepsilon_t,
  \label{eq:arma}
\end{equation}
with $\Phi$ and $\Theta$ sharing no common roots.
\end{definition}

The ARMA($p,q$) ACF decays geometrically for lags $h > q - p$ (when
$q \geq p$), and the PACF decays geometrically for lags $h > p - q$
(when $p \geq q$). This \emph{mixing} of tail behaviors is what
distinguishes ARMA from pure AR or MA, and motivates information-criteria
model selection rather than sharp cutoff identification.

The Wold MA representation coefficients satisfy:
\begin{equation}
  \psi(z) = \frac{\Theta(z)}{\Phi(z)} = \sum_{j=0}^{\infty}\psi_j z^j.
  \label{eq:wold_coeffs}
\end{equation}

\subsection{Integrated and Seasonal Models}

Many economic and business series exhibit \emph{stochastic trends}: the first
difference $\nabla Y_t = Y_t - Y_{t-1}$ is stationary even though $Y_t$ itself
is not.

\begin{definition}[ARIMA($p,d,q$) Process]
$\{Y_t\}$ follows an ARIMA($p,d,q$) process if $\nabla^d Y_t$ is a
covariance-stationary ARMA($p,q$) process:
\begin{equation}
  \Phi(B)\nabla^d Y_t = c + \Theta(B)\varepsilon_t.
  \label{eq:arima}
\end{equation}
\end{definition}

The special case ARIMA(0,1,0), i.e., $Y_t = Y_{t-1} + \varepsilon_t$, is the
\emph{random walk}, the canonical model for efficient-market asset prices.

Seasonal patterns with period $s$ are accommodated by the
\emph{SARIMA($p,d,q$)$\times$($P,D,Q$)$_s$} model~\cite{box2015time}:
\begin{equation}
  \Phi(B)\Phi_s(B^s)\nabla^d\nabla_s^D Y_t
  = \Theta(B)\Theta_s(B^s)\varepsilon_t,
  \label{eq:sarima}
\end{equation}
where $\nabla_s = 1 - B^s$, $\Phi_s(z) = 1 - \Phi_1 z - \cdots - \Phi_P z^P$,
and $\Theta_s$ is defined analogously. For monthly data, $s = 12$; for
quarterly data, $s = 4$.

\begin{definition}[Unit Root]
A process has a \emph{unit root} if $\Phi(z) = 0$ has a root at $z = 1$,
equivalently the characteristic polynomial evaluated at 1 equals zero. Unit
roots produce non-stationarity manifested as a stochastic trend.
\end{definition}

The \emph{Augmented Dickey--Fuller} (ADF) test~\cite{dickey1979distribution}
evaluates $H_0$: unit root against $H_1$: stationarity via the regression
\begin{equation}
  \nabla Y_t = \alpha + \beta t + \delta Y_{t-1}
               + \sum_{j=1}^{k}\gamma_j \nabla Y_{t-j} + u_t,
  \label{eq:adf}
\end{equation}
with test statistic $\tau = \hat{\delta}/\mathrm{se}(\hat{\delta})$ compared
to Dickey--Fuller critical values.

\subsubsection{Co-integration and Error-Correction Models}

When two $I(1)$ series $\{X_t\}$ and $\{Y_t\}$ share a long-run equilibrium
$Y_t - \beta X_t \sim I(0)$, they are \emph{co-integrated}~\cite{engle1987co}.
The \emph{vector error-correction model} (VECM) is:
\begin{equation}
  \nabla\mathbf{z}_t
  = \alpha\boldsymbol{\beta}^{\top}\mathbf{z}_{t-1}
  + \sum_{j=1}^{k-1}\Gamma_j\nabla\mathbf{z}_{t-j} + \boldsymbol{\varepsilon}_t,
  \label{eq:vecm}
\end{equation}
where $\mathbf{z}_t = (X_t, Y_t)^{\top}$, $\boldsymbol{\beta}$ is the
co-integrating vector, and $\alpha$ is the adjustment (loading) matrix.
The $\boldsymbol{\beta}^{\top}\mathbf{z}_{t-1}$ term is the
\emph{error-correction} component restoring equilibrium after shocks.

%=============================================================================
\section{Algorithms and Methods}
\label{sec:algorithms}

\subsection{Model Identification}

The Box--Jenkins identification procedure uses the sample ACF and PACF
(Table~\ref{tab:acf_patterns}) together with information criteria.

\begin{table}[htbp]
\centering
\caption{ACF and PACF Signature Patterns for Model Identification}
\label{tab:acf_patterns}
\begin{tabular}{lccc}
\toprule
\textbf{Model} & \textbf{ACF} & \textbf{PACF} & \textbf{Characteristic}\\
\midrule
AR($p$)    & Tails off (geometric) & Cuts off after $p$ & $p$ spikes in PACF \\
MA($q$)    & Cuts off after $q$    & Tails off          & $q$ spikes in ACF  \\
ARMA($p,q$)& Tails off (lag $>q$)  & Tails off (lag $>p$) & Both tail off     \\
White noise& No significant lags   & No significant lags & Flat              \\
Random walk& Slow decay            & Spike at lag 1     & Non-stationary    \\
\bottomrule
\end{tabular}
\end{table}

\subsection{Algorithm 1: Yule--Walker Estimation for AR($p$)}

\noindent\textbf{Input:} Observed series $y_1,\ldots,y_T$; order $p$.\\
\noindent\textbf{Output:} Parameter estimates $\hat{\phi}_1,\ldots,\hat{\phi}_p$,
$\hat{\sigma}^2$.

\begin{enumerate}
  \item Compute sample autocovariances
        $\hat{\gamma}(h) = T^{-1}\sum_{t=h+1}^T (y_t-\bar{y})(y_{t-h}-\bar{y})$
        for $h = 0, 1, \ldots, p$.
  \item Form the sample autocorrelation matrix
        $\hat{\mathbf{R}} \in \mathbb{R}^{p\times p}$ with
        $\hat{R}_{ij} = \hat{\rho}(|i-j|)$.
  \item Solve the \emph{Yule--Walker system}
        \begin{equation}
          \hat{\mathbf{R}}\,\hat{\boldsymbol{\phi}}
          = \hat{\boldsymbol{\rho}},
          \label{eq:yw_matrix}
        \end{equation}
        where $\hat{\boldsymbol{\rho}} = (\hat{\rho}(1),\ldots,\hat{\rho}(p))^{\top}$,
        via Cholesky decomposition or the Levinson--Durbin recursion~\cite{durbin1960fitting}.
  \item Estimate the noise variance:
        $\hat{\sigma}^2 = \hat{\gamma}(0)(1 - \hat{\boldsymbol{\phi}}^{\top}
        \hat{\boldsymbol{\rho}})$.
\end{enumerate}

The Yule--Walker estimator is $\sqrt{T}$-consistent under stationarity.

\subsection{Algorithm 2: Maximum Likelihood Estimation for ARMA($p,q$)}

\noindent\textbf{Input:} Series $\mathbf{y}$; orders $p,q$;
Gaussian innovations assumed.\\
\noindent\textbf{Output:} $(\hat{\boldsymbol{\phi}}, \hat{\boldsymbol{\theta}}, \hat{\sigma}^2)$.

\begin{enumerate}
  \item Initialize the innovation recursion. Set $\hat{\varepsilon}_1 = \cdots
        = \hat{\varepsilon}_{-\max(p,q)+1} = 0$.
  \item For $t = 1, \ldots, T$, compute one-step-ahead prediction errors:
        \begin{equation}
          \hat{\varepsilon}_t = y_t - c - \sum_{i=1}^p \phi_i y_{t-i}
                                - \sum_{j=1}^q \theta_j \hat{\varepsilon}_{t-j}.
          \label{eq:innovations}
        \end{equation}
  \item Compute the conditional log-likelihood:
        \begin{equation}
          \ell(\boldsymbol{\phi},\boldsymbol{\theta},\sigma^2 \mid \mathbf{y})
          = -\frac{T}{2}\log(2\pi\sigma^2)
            - \frac{1}{2\sigma^2}\sum_{t=1}^T \hat{\varepsilon}_t^2.
          \label{eq:cll}
        \end{equation}
  \item Maximize \eqref{eq:cll} over $(\boldsymbol{\phi},\boldsymbol{\theta})$
        using a quasi-Newton method (e.g., L-BFGS) with stationarity and
        invertibility constraints enforced via reparameterization
        (e.g., partial autocorrelation space for AR coefficients~\cite{jones1980maximum}).
  \item Recover $\hat{\sigma}^2 = T^{-1}\sum_t \hat{\varepsilon}_t^2$ at
        the optimum.
\end{enumerate}

\textbf{Exact MLE} replaces step 3 with the multivariate Gaussian likelihood
$\ell = -\frac{1}{2}\bigl[\log\det\boldsymbol{\Sigma} +
(\mathbf{y}-\boldsymbol{\mu})^{\top}\boldsymbol{\Sigma}^{-1}(\mathbf{y}-\boldsymbol{\mu})\bigr]$,
where $\boldsymbol{\Sigma}$ is the Toeplitz covariance matrix;
the Durbin--Levinson algorithm computes this in $O(T)$ time~\cite{brockwell1991time}.

\subsection{Algorithm 3: Model Selection via Information Criteria}

Two primary criteria balance fit and parsimony:
\begin{align}
  \mathrm{AIC}(p,q) &= -2\hat{\ell} + 2(p+q+1), \label{eq:aic}\\
  \mathrm{BIC}(p,q) &= -2\hat{\ell} + (p+q+1)\log T. \label{eq:bic}
\end{align}
The AIC~\cite{akaike1974new} targets one-step-ahead prediction accuracy
asymptotically, while the BIC~\cite{schwarz1978estimating} is consistent
for the true order. In practice, a grid search over $p \leq p_{\max}$,
$q \leq q_{\max}$ is performed, selecting the model minimizing AIC or BIC.

\subsection{Algorithm 4: Multi-Step Forecasting}

Given parameter estimates and data through time $T$, the MMSE
$h$-step-ahead forecast $\hat{y}_{T+h|T} = \mathbb{E}[Y_{T+h}\mid y_1,\ldots,y_T]$
is obtained recursively:

\begin{enumerate}
  \item Set $\hat{y}_{T+j|T} = y_{T+j}$ for $j \leq 0$ and
        $\hat{\varepsilon}_{T+j|T} = 0$ for $j > 0$.
  \item For $h = 1, 2, \ldots, H$, compute
        \begin{equation}
          \hat{y}_{T+h|T} = \hat{c} + \sum_{i=1}^p \hat{\phi}_i\,\hat{y}_{T+h-i|T}
                            + \sum_{j=1}^q \hat{\theta}_j\,\hat{\varepsilon}_{T+h-j|T}.
          \label{eq:forecast_recursion}
        \end{equation}
  \item Compute forecast error variances
        $v_h = \sigma^2\sum_{j=0}^{h-1}\psi_j^2$ using Wold coefficients
        $\{\psi_j\}$ from \eqref{eq:wold_coeffs}.
  \item Form $(1-\alpha)$ prediction intervals:
        $\hat{y}_{T+h|T} \pm z_{\alpha/2}\sqrt{v_h}$.
\end{enumerate}

For ARIMA models, forecasts are computed on the differenced series and then
\emph{integrated} (cumulative-summed) back to the original scale.

\subsection{Algorithm 5: Holt--Winters Exponential Smoothing}

Exponential smoothing~\cite{holt2004forecasting,winters1960forecasting} provides
a robust, computation-friendly alternative to ARIMA. The additive
Holt--Winters model with trend and seasonality (period $s$) maintains:
\begin{align}
  \ell_t &= \alpha(y_t - s_{t-m}) + (1-\alpha)(\ell_{t-1} + b_{t-1}),
            \label{eq:hw_level}\\
  b_t    &= \beta^*(\ell_t - \ell_{t-1}) + (1-\beta^*)b_{t-1},
            \label{eq:hw_trend}\\
  s_t    &= \gamma(y_t - \ell_{t-1} - b_{t-1}) + (1-\gamma)s_{t-m},
            \label{eq:hw_seasonal}
\end{align}
with $h$-step forecast
$\hat{y}_{T+h|T} = \ell_T + hb_T + s_{T+h-m\lceil h/m\rceil}$.
Parameters $\alpha, \beta^*, \gamma \in (0,1)$ are estimated by minimizing
$\sum_t(\hat{y}_{t|t-1} - y_t)^2$.

Hyland and Gardner~\cite{gardner1985exponential} and Hyndman
et al.~\cite{hyndman2008forecasting} showed that every exponential smoothing
method corresponds to a state-space model with a specific error structure
(additive or multiplicative), unifying smoothing with ARIMA under the
ETS (Error, Trend, Seasonality) framework.

%=============================================================================
\section{Worked Examples}
\label{sec:examples}

\subsection{Example 1: AR(2) Parameter Estimation}

Consider the zero-mean AR(2) model $Y_t = 0.6Y_{t-1} - 0.3Y_{t-2} + \varepsilon_t$
with $\sigma^2 = 1$.

\paragraph{Stationarity check.}
The AR polynomial is $\Phi(z) = 1 - 0.6z + 0.3z^2$. Setting $\Phi(z) = 0$:
\begin{align*}
  z &= \frac{0.6 \pm \sqrt{0.36 - 1.2}}{0.6}
     = \frac{0.6 \pm \sqrt{-0.84}}{0.6}.
\end{align*}
The roots are complex conjugates with modulus
$|z| = 1/\sqrt{0.3} \approx 1.826 > 1$,
confirming stationarity.

\paragraph{Yule--Walker solution.}
From \eqref{eq:yw}:
\begin{align*}
  \rho(1) &= \phi_1 + \phi_2\rho(1) \;\Rightarrow\; \rho(1) = \frac{0.6}{1+0.3} \approx 0.462,\\
  \rho(2) &= \phi_1\rho(1) + \phi_2 = 0.6(0.462) - 0.3 \approx -0.023.
\end{align*}

\paragraph{Theoretical PACF.}
By definition, $\phi_{11} = \rho(1) \approx 0.462$,
$\phi_{22} = \phi_2 = -0.3$, and $\phi_{hh} = 0$ for $h \geq 3$,
consistent with the AR(2) signature in Table~\ref{tab:acf_patterns}.

\paragraph{Variance.}
$\gamma(0) = \sigma^2/(1 - \phi_1\rho(1) - \phi_2\rho(2)) \approx 1/(1 - 0.277 + 0.007) \approx 1.36$.

\subsection{Example 2: ARIMA Forecasting for Airline Passenger Data}

The \texttt{AirPassengers} dataset (monthly international airline passengers,
1949--1960, $T=144$) is a canonical benchmark~\cite{box2015time}.
The series exhibits an exponential trend and multiplicative seasonality
(variance increasing with level), motivating a log transformation.
Let $x_t = \log(y_t)$.

\paragraph{Identification.}
The ACF of $x_t$ decays very slowly (unit-root behavior). After
first-differencing ($d=1$), the seasonal ACF at lag 12 remains significant,
motivating a seasonal difference ($D=1$, $s=12$). On $w_t = \nabla\nabla_{12}x_t$:
the ACF has significant spikes at lags 1 and 12; the PACF decays at
both non-seasonal and seasonal lags. This pattern suggests SARIMA(0,1,1)$\times$(0,1,1)$_{12}$,
the \emph{airline model}~\cite{box2015time}.

\paragraph{Estimated model.}
Conditional MLE yields (standard errors in parentheses):
$\hat{\theta}_1 = -0.402\;(0.089)$, $\hat{\Theta}_1 = -0.557\;(0.073)$,
$\hat{\sigma}^2 = 0.00134$.

\paragraph{Forecasts.}
Twelve-month-ahead point forecasts on the log scale are back-transformed via
$\hat{y}_{T+h|T} = \exp(\hat{x}_{T+h|T} + \hat{\sigma}^2 v_h / 2)$
(bias correction for lognormal distribution). The 95\% prediction bands widen
appropriately with horizon.

\paragraph{Diagnostics.}
Ljung--Box test~\cite{ljung1978measure}:
\begin{equation}
  Q(m) = T(T+2)\sum_{h=1}^m \frac{\hat{\rho}_\varepsilon^2(h)}{T-h},
  \label{eq:ljungbox}
\end{equation}
yields $Q(20) = 17.4$ with $p = 0.62$ under $\chi^2_{18}$, confirming
adequate fit.

\subsection{Example 3: Selecting Between Models Using AIC/BIC}

Table~\ref{tab:aic_bic} shows AIC and BIC values for candidate ARMA models
fitted to a simulated ARMA(1,1) series of length $T=200$.

\begin{table}[htbp]
\centering
\caption{AIC and BIC for Candidate ARMA Models ($T=200$, True Model: ARMA(1,1))}
\label{tab:aic_bic}
\begin{tabular}{ccccc}
\toprule
\textbf{Model} & \textbf{$p+q$} & \textbf{Log-lik} & \textbf{AIC} & \textbf{BIC}\\
\midrule
AR(1)         & 1 & $-278.4$ & 560.8 & 567.2 \\
MA(1)         & 1 & $-279.1$ & 562.2 & 568.6 \\
ARMA(1,1)     & 2 & $-271.3$ & 548.6 & 558.2 \\
ARMA(2,1)     & 3 & $-271.0$ & 550.0 & 562.8 \\
ARMA(1,2)     & 3 & $-271.2$ & 550.4 & 563.2 \\
ARMA(2,2)     & 4 & $-270.8$ & 551.6 & 567.6 \\
\bottomrule
\end{tabular}
\end{table}

Both AIC and BIC correctly select ARMA(1,1). Note that AIC penalizes less
heavily and would be more likely to select a higher-order model in borderline
cases; BIC's stronger penalty is consistent with order consistency.

%=============================================================================
\section{Applications in Data Science}
\label{sec:applications}

\subsection{Demand and Supply Chain Forecasting}

Retail demand series typically exhibit weekly or annual seasonality and
trend. SARIMA and Holt--Winters models are workhorses for point-of-sale
inventory management. Makridakis et al.~\cite{makridakis2018m4} evaluated
over 60,000 time series in the M4 competition, finding that exponential
smoothing methods remain highly competitive, often outperforming complex
machine-learning approaches on short to medium horizons.

\subsection{Financial Econometrics}

Asset log-returns are approximately white noise (efficient market hypothesis)
but exhibit volatility clustering motivating GARCH extensions. However,
ARIMA models remain essential for: forecasting macroeconomic indicators
(GDP, CPI, unemployment) that exhibit autocorrelation, modeling the term
structure of interest rates, and detecting mean-reversion. The random walk
model $d = 1$ is the standard benchmark against which more complex financial
models are evaluated~\cite{hamilton1994time}.

\subsection{Epidemiological Surveillance}

Time-series models are deployed in public health to detect anomalous
disease counts above seasonal baselines. The Farrington
algorithm~\cite{farrington1996statistical} adapts a quasi-Poisson regression
with ARIMA residuals to detect outbreak signals in syndromic surveillance
data. During the COVID-19 pandemic, ARIMA and exponential smoothing models
were used as rapid-deployment baselines for excess mortality estimation.

\subsection{Environmental Monitoring and Climate Science}

Monthly mean temperature, sea-level, and atmospheric CO$_2$ series
(e.g., Mauna Loa keeling curve) are modeled as ARIMA processes with
deterministic seasonal harmonics. Unit-root tests determine whether
long-term changes represent stochastic trends (requiring differencing) or
deterministic trends (requiring time-trend regression). Co-integration
analysis is used to study long-run equilibria between climate variables such
as global temperature and CO$_2$ concentration~\cite{engle1987co}.

\subsection{Industrial Process Control and IoT}

Manufacturing sensor streams (vibration, temperature, pressure) are
monitored using ARIMA-based residual control charts. An out-of-control
signal is flagged when the one-step-ahead prediction error exceeds control
limits derived from the fitted model variance. The \emph{transfer function
model}~\cite{box2015time}, an extension of ARIMA with exogenous inputs,
enables modeling of cause-and-effect relationships between process inputs
and outputs.

%=============================================================================
\section{Limitations and Assumptions}
\label{sec:limits}

\subsection{Linearity}

ARMA-class models impose a \emph{linear} conditional mean structure.
Nonlinear dynamics such as threshold behavior (business cycle asymmetry),
regime switching (financial crises), or chaotic systems cannot be captured
without extensions such as TAR~\cite{tong1990non},
SETAR, or Markov-switching models.

\subsection{Gaussianity}

Standard inference and exact likelihood computation assume Gaussian
innovations. Financial returns, count data, and heavy-tailed sensor readings
violate this assumption. While QML (quasi-maximum likelihood) remains
consistent under misspecification of the innovation distribution, prediction
intervals based on Gaussian quantiles are unreliable for heavy-tailed series.

\subsection{Constant Variance (Homoscedasticity)}

ARIMA models assume $\mathrm{Var}(\varepsilon_t) = \sigma^2$ constant over
time. Volatility clustering in financial data violates this, necessitating
ARCH/GARCH models~\cite{engle1982autoregressive,bollerslev1986generalized}.

\subsection{Stationarity After Differencing}

The ARIMA framework assumes that a suitable integer number of differences $d$
renders the series stationary. Fractionally integrated processes (long memory)
require ARFIMA models with non-integer $d$~\cite{granger1980introduction}.
Additionally, structural breaks in mean or variance can be misidentified as
unit roots~\cite{perron1989great}.

\subsection{Parameter Stability}

The model assumes \emph{time-invariant} parameters. In practice, economic
regimes shift, technologies change, and consumer behavior evolves. Rolling
estimation, recursive forecasting, and state-space models with time-varying
parameters are used to address parameter instability~\cite{durbin2012time}.

\subsection{Univariate versus Multivariate}

Classical ARIMA treats series one at a time, ignoring information in related
series. Vector autoregressive (VAR) models extend AR to multivariate settings,
capturing Granger causality. The VECM extends VAR to co-integrated systems.

\subsection{Sample Size and Overfitting}

High-order ARMA models have many parameters; with $p + q + 1$ free parameters
in the mean equation alone, overfitting is a concern for short series.
BIC-based selection and cross-validation (e.g., time-series split) mitigate
overfitting.

%=============================================================================
\section{Connections to Related Methods}
\label{sec:related}

\subsection{State-Space and Structural Time-Series Models}

Every ARIMA model admits a state-space representation
\begin{align}
  \boldsymbol{\alpha}_{t+1} &= \mathbf{T}\boldsymbol{\alpha}_t
                              + \mathbf{R}\boldsymbol{\eta}_t, \label{eq:ss_trans}\\
  y_t &= \mathbf{Z}^{\top}\boldsymbol{\alpha}_t + \varepsilon_t, \label{eq:ss_obs}
\end{align}
enabling \emph{Kalman filter} recursions for exact likelihood evaluation
and optimal state estimation~\cite{kalman1960new,durbin2012time}.
Harvey's structural time-series
models~\cite{harvey1990forecasting} specify trend, seasonal, and
irregular components directly in \eqref{eq:ss_trans}--\eqref{eq:ss_obs},
providing an interpretable decomposition and accommodating missing data
and irregular spacing naturally.

\subsection{ARCH and GARCH Models}

Engle's~\cite{engle1982autoregressive} ARCH model specifies conditional
heteroscedasticity:
$\varepsilon_t = \sigma_t z_t$ with $z_t \stackrel{\text{iid}}{\sim}(0,1)$ and
$\sigma_t^2 = \omega + \sum_{i=1}^r \alpha_i \varepsilon_{t-i}^2$.
Bollerslev's~\cite{bollerslev1986generalized} GARCH($r,s$) adds
$\sum_{j=1}^s \beta_j \sigma_{t-j}^2$.
ARMA-GARCH models combine autoregressive mean dynamics with GARCH variance
dynamics, the standard model for financial log-return series.

\subsection{Vector Autoregression (VAR)}

The VAR($p$) model for $K$-dimensional $\mathbf{y}_t$:
\begin{equation}
  \mathbf{y}_t = \mathbf{c} + \sum_{i=1}^p \mathbf{A}_i \mathbf{y}_{t-i}
                 + \boldsymbol{\varepsilon}_t
  \label{eq:var}
\end{equation}
generalizes univariate AR to capture cross-series dynamics.
VAR is the standard tool for structural macroeconomic analysis via impulse
response functions and forecast error variance decompositions~\cite{sims1980macroeconomics,lutkepohl2005new}.

\subsection{Machine Learning Approaches}

Modern approaches apply recurrent neural networks
(RNN/LSTM~\cite{hochreiter1997long}) and Transformer architectures to
time-series, capturing nonlinear and long-range dependencies. However,
classical models retain advantages in interpretability, statistical
efficiency for small samples, and robustness. Hybrid approaches combine
classical decomposition with neural networks for residuals (e.g., N-BEATS,
N-HiTS~\cite{oreshkin2019nbeats}). The M4 and M5 competitions showed that
ensembles of classical and neural models consistently outperform either
family alone~\cite{makridakis2018m4}.

\subsection{Spectral Analysis}

The \emph{spectral density} of a covariance-stationary process is
$f(\omega) = \frac{1}{2\pi}\sum_{h=-\infty}^{\infty}\gamma(h)e^{-i\omega h}$,
the Fourier transform of the ACVF. For an ARMA($p,q$) process:
\begin{equation}
  f(\omega) = \frac{\sigma^2}{2\pi}
              \frac{|\Theta(e^{-i\omega})|^2}{|\Phi(e^{-i\omega})|^2}.
  \label{eq:spectral}
\end{equation}
Spectral analysis identifies periodic components and provides a
complementary view to the time-domain ACF/PACF representation.

%=============================================================================
\section{Conclusion}
\label{sec:conclusion}

This survey has presented a unified, mathematically rigorous treatment of
classical time-series models from the AR, MA, and ARMA families through
their integrated and seasonal extensions (ARIMA/SARIMA), exponential
smoothing, and co-integration/error-correction models.

We have shown how these models address the four core themes of predictive
modeling for dependent data. First, the ACF, PACF, and unit-root tests
(ADF) provide principled tools for \emph{identifying dependencies and
non-stationarity} in structured temporal data. Second, the ARIMA and
Holt--Winters frameworks offer a rich menu of \emph{appropriate model forms}
matched to trend, seasonality, and autocorrelation patterns. Third, the MMSE
forecast recursion and state-space formulation yield \emph{optimal forecasts
with calibrated prediction intervals} that properly account for serial
dependence. Fourth, the state-space and structural decomposition
representations \emph{reveal latent trends, seasonal cycles, and irregular
components} in otherwise complex series.

Despite the rise of deep learning for time-series, classical models remain
indispensable: they provide interpretable parameterizations, asymptotically
valid inference, and competitive forecast accuracy particularly on shorter
horizons or smaller datasets. Their assumptions --- linearity, homoscedasticity,
and Gaussian innovations --- are well understood and motivate principled
extensions to GARCH, state-space, and machine-learning hybrids.

For practitioners in data science, the Box--Jenkins identification--estimation--%
diagnostic cycle remains a valuable discipline: it forces explicit
consideration of data properties, model adequacy, and forecast uncertainty
before deployment.

%=============================================================================
\bibliographystyle{IEEEtran}
\begin{thebibliography}{99}

\bibitem{akaike1974new}
H. Akaike, ``A new look at the statistical model identification,''
\textit{IEEE Transactions on Automatic Control}, vol.~19, no.~6,
pp.~716--723, 1974.

\bibitem{bollerslev1986generalized}
T. Bollerslev, ``Generalized autoregressive conditional heteroskedasticity,''
\textit{Journal of Econometrics}, vol.~31, no.~3, pp.~307--327, 1986.

\bibitem{box2015time}
G. E. P. Box, G. M. Jenkins, G. C. Reinsel, and G. M. Ljung,
\textit{Time Series Analysis: Forecasting and Control}, 5th ed.
Hoboken, NJ, USA: Wiley, 2015.

\bibitem{brockwell1991time}
P. J. Brockwell and R. A. Davis,
\textit{Time Series: Theory and Methods}, 2nd ed.
New York, NY, USA: Springer, 1991.

\bibitem{dickey1979distribution}
D. A. Dickey and W. A. Fuller, ``Distribution of the estimators for
autoregressive time series with a unit root,''
\textit{Journal of the American Statistical Association}, vol.~74,
no.~366, pp.~427--431, 1979.

\bibitem{durbin1960fitting}
J. Durbin, ``The fitting of time-series models,''
\textit{Review of the International Statistical Institute}, vol.~28,
no.~3, pp.~233--244, 1960.

\bibitem{durbin2012time}
J. Durbin and S. J. Koopman,
\textit{Time Series Analysis by State Space Methods}, 2nd ed.
Oxford, UK: Oxford University Press, 2012.

\bibitem{engle1982autoregressive}
R. F. Engle, ``Autoregressive conditional heteroscedasticity with estimates
of the variance of United Kingdom inflation,''
\textit{Econometrica}, vol.~50, no.~4, pp.~987--1007, 1982.

\bibitem{engle1987co}
R. F. Engle and C. W. J. Granger, ``Co-integration and error correction:
representation, estimation, and testing,''
\textit{Econometrica}, vol.~55, no.~2, pp.~251--276, 1987.

\bibitem{farrington1996statistical}
C. P. Farrington, N. Andrews, A. D. Beale, and M. A. Catchpole,
``A statistical algorithm for the early detection of outbreaks of
infectious disease,'' \textit{Journal of the Royal Statistical Society:
Series A}, vol.~159, no.~3, pp.~547--563, 1996.

\bibitem{gardner1985exponential}
E. S. Gardner, ``Exponential smoothing: the state of the art,''
\textit{Journal of Forecasting}, vol.~4, no.~1, pp.~1--28, 1985.

\bibitem{granger1980introduction}
C. W. J. Granger and R. Joyeux, ``An introduction to long-memory time
series models and fractional differencing,''
\textit{Journal of Time Series Analysis}, vol.~1, no.~1, pp.~15--29, 1980.

\bibitem{hamilton1994time}
J. D. Hamilton,
\textit{Time Series Analysis}.
Princeton, NJ, USA: Princeton University Press, 1994.

\bibitem{harvey1990forecasting}
A. C. Harvey,
\textit{Forecasting, Structural Time Series Models and the Kalman Filter}.
Cambridge, UK: Cambridge University Press, 1990.

\bibitem{hochreiter1997long}
S. Hochreiter and J. Schmidhuber, ``Long short-term memory,''
\textit{Neural Computation}, vol.~9, no.~8, pp.~1735--1780, 1997.

\bibitem{holt2004forecasting}
C. C. Holt, ``Forecasting seasonals and trends by exponentially weighted
moving averages,'' \textit{International Journal of Forecasting}, vol.~20,
no.~1, pp.~5--10, 2004.

\bibitem{hyndman2008forecasting}
R. J. Hyndman, A. B. Koehler, J. K. Ord, and R. D. Snyder,
\textit{Forecasting with Exponential Smoothing: The State Space Approach}.
Berlin, Germany: Springer, 2008.

\bibitem{jones1980maximum}
R. H. Jones, ``Maximum likelihood fitting of ARMA models to time series
with missing observations,''
\textit{Technometrics}, vol.~22, no.~3, pp.~389--395, 1980.

\bibitem{kalman1960new}
R. E. Kalman, ``A new approach to linear filtering and prediction problems,''
\textit{Journal of Basic Engineering}, vol.~82, no.~1, pp.~35--45, 1960.

\bibitem{ljung1978measure}
G. M. Ljung and G. E. P. Box, ``On a measure of lack of fit in time series
models,'' \textit{Biometrika}, vol.~65, no.~2, pp.~297--303, 1978.

\bibitem{lutkepohl2005new}
H. L{\"u}tkepohl,
\textit{New Introduction to Multiple Time Series Analysis}.
Berlin, Germany: Springer, 2005.

\bibitem{makridakis2018m4}
S. Makridakis, E. Spiliotis, and V. Assimakopoulos, ``The M4 competition:
results, findings, conclusion and way forward,''
\textit{International Journal of Forecasting}, vol.~34, no.~4,
pp.~802--808, 2018.

\bibitem{oreshkin2019nbeats}
B. N. Oreshkin, D. Carpov, N. Chapados, and Y. Bengio,
``N-BEATS: neural basis expansion analysis for interpretable time series
forecasting,'' in \textit{Proc. 8th International Conference on Learning
Representations (ICLR)}, 2020.

\bibitem{perron1989great}
P. Perron, ``The great crash, the oil price shock, and the unit root
hypothesis,'' \textit{Econometrica}, vol.~57, no.~6, pp.~1361--1401, 1989.

\bibitem{schuster1898investigation}
A. Schuster, ``On the investigation of hidden periodicities with application
to a supposed 26-day period of meteorological phenomena,''
\textit{Terrestrial Magnetism}, vol.~3, no.~1, pp.~13--41, 1898.

\bibitem{schwarz1978estimating}
G. Schwarz, ``Estimating the dimension of a model,''
\textit{The Annals of Statistics}, vol.~6, no.~2, pp.~461--464, 1978.

\bibitem{sims1980macroeconomics}
C. A. Sims, ``Macroeconomics and reality,''
\textit{Econometrica}, vol.~48, no.~1, pp.~1--48, 1980.

\bibitem{slutzky1927summation}
E. Slutzky, ``The summation of random causes as the source of cyclic
processes,'' \textit{Econometrica}, vol.~5, no.~2, pp.~105--146, 1937.

\bibitem{tong1990non}
H. Tong,
\textit{Non-linear Time Series: A Dynamical System Approach}.
Oxford, UK: Oxford University Press, 1990.

\bibitem{winters1960forecasting}
P. R. Winters, ``Forecasting sales by exponentially weighted moving averages,''
\textit{Management Science}, vol.~6, no.~3, pp.~324--342, 1960.

\bibitem{wold1938study}
H. Wold,
\textit{A Study in the Analysis of Stationary Time Series}.
Stockholm, Sweden: Almqvist \& Wiksell, 1938.

\bibitem{yule1927periodicity}
G. U. Yule, ``On a method of investigating periodicities in disturbed
series, with special reference to Wolfer's sunspot numbers,''
\textit{Philosophical Transactions of the Royal Society of London A},
vol.~226, pp.~267--298, 1927.

\end{thebibliography}

\end{document}
